\section{Conventions and endpoint normalization}\label{shj:conventions}
Throughout,
\[
0<a<1,\qquad b=1-a,\qquad \ell=\log(2\pi),\qquad
\kappa=\gamma+\ell,\qquad B_2(a)=a^2-a+\frac16.
\]
The fractional part $\{x+a\}$ is interpreted periodically; its value at the single jump point does not affect ordinary integrals. The notation
\[
\zeta^{[j]}(s,x)=\partial_s^j\zeta(s,x),\qquad
\Li_s^{[j]}(z)=\partial_s^j\Li_s(z)
\]
always means \emph{spectral} differentiation. In contrast, $\gamma_m^{(r)}(x)$ and $\psi^{(r)}(x)$ mean derivatives in the real argument. Euler's constant is $\gamma=\gamma_0(1)$, not a new parameter.

The generalized Stieltjes constants are normalized by
\begin{equation}\label{shj:laurent}
\zeta(1+u,x)=\frac1u+R(u,x),\qquad
R(u,x)=\sum_{m\ge0}\frac{(-1)^m}{m!}\gamma_m(x)u^m.
\end{equation}
Thus $\gamma_0(x)=-\psi(x)$. The Hurwitz shift identity gives the exact local decomposition
\begin{equation}\label{shj:localgamma}
\gamma_m(x)=\frac{\log^m x}{x}+\gamma_m(1+x),\qquad x>0.
\end{equation}
Every logarithm of a positive real argument is the real logarithm. For $z=e^{2\pi\iu a}\ne1$, $\Li_s(z)$ and $\Li_s(z^{-1})$ denote their usual boundary values from the unit disk, continued in $s$; these are entire in $s$. For real $s$ they are conjugate, but we retain them separately when $s$ is complex.

\begin{definition}[Canonical shifted Stieltjes finite part]\label{shj:defJ}
For nonnegative integers $m,n$, define
\begin{align}
J_{mn}(a)=\lim_{\epsilon\downarrow0}\bigg\{&
\int_\epsilon^b\gamma_m(x)\gamma_n(x+a)\dd x
+\int_\epsilon^a\gamma_m(b+t)\gamma_n(t)\dd t\notag\\
&+\frac{\gamma_n(a)}{m+1}\log^{m+1}\epsilon
+\frac{\gamma_m(b)}{n+1}\log^{n+1}\epsilon\bigg\}.
\label{shj:Jcutoff}
\end{align}
Equivalently, this is denoted
$\FP\int_0^1\gamma_m(x)\gamma_n(\{x+a\})\dd x$.
\end{definition}
The two cutoffs can tend to zero independently: each counterterm belongs to its own local endpoint. The coordinates are exactly $x$ and $t=x-b$, with no rescaling or additional subtraction length. A different finite-part convention can change local constants and therefore change the identities.

\begin{lemma}[Existence and an ordinary-integral version]\label{shj:Jexists}
The limit \eqref{shj:Jcutoff} exists and is equivalent to
\begin{align}
J_{mn}(a)={}&\int_0^b\left[
\gamma_m(x)\gamma_n(x+a)-\gamma_n(a)\frac{\log^m x}{x}\right]\dd x\notag\\
&+\int_0^a\left[
\gamma_m(b+t)\gamma_n(t)-\gamma_m(b)\frac{\log^n t}{t}\right]\dd t\notag\\
&+\frac{\gamma_n(a)}{m+1}\log^{m+1}b
+\frac{\gamma_m(b)}{n+1}\log^{n+1}a.
\label{shj:Jregular}
\end{align}
Both integrals on the right are ordinary improper integrals. Moreover,
$J_{mn}(a)=J_{nm}(1-a)$.
\end{lemma}
\begin{proof}
By \eqref{shj:localgamma}, subtracting the singular term in the first integral leaves $O(1+|\log x|^m)$ near zero: the other factor differs from $\gamma_n(a)$ by $O(x)$. The second integral has the same property. Integrating $\log^m x/x$ explicitly yields \eqref{shj:Jregular}. Interchanging the two pieces and replacing $a$ by $b$ proves the symmetry.
\end{proof}
We use analogous finite parts for polygamma and derivative products: expand each smooth factor at the distinct singular point, integrate the resulting powers and logarithms, and retain the constant term after deleting all divergent cutoff terms. Section~\ref{stconv:sec:contact} makes this prescription precise through periodic distributions.

\section{The shifted Hurwitz correlation}\label{shj:master}
\begin{proposition}[Classical Fourier method, shifted form]\label{shj:masterthm}
For $\Re s<1$, $\Re t<1$, define the ordinary integral
\begin{equation}\label{shj:Cdef}
C(s,t;a)=\int_0^1\zeta(s,x)\zeta(t,\{x+a\})\dd x.
\end{equation}
It is jointly holomorphic in this domain. Put $w=2-s-t$ and $z=e^{2\pi\iu a}$. Then
\begin{align}
C(s,t;a)=\frac{\Gamma(1-s)\Gamma(1-t)}{(2\pi)^w}
\bigg[&e^{-\pi\iu(s-t)/2}\Li_w(z)\notag\\
&+e^{\pi\iu(s-t)/2}\Li_w(z^{-1})\bigg].
\label{shj:Cpoly}
\end{align}
It also has the meromorphic representation
\begin{align}
C(s,t;a)={}&B(1-s,s+t-1)\zeta(s+t-1,a)\notag\\
&+B(1-t,s+t-1)\zeta(s+t-1,b).
\label{shj:Cbeta}
\end{align}
At exceptional arguments of the beta factors, the \emph{sum} in \eqref{shj:Cbeta} is continued before evaluation. Apparent poles of separate summands need not be poles of $C$.
\end{proposition}
\begin{proof}
Split the integral at $b$. At $x=0$ only the first factor is singular and has the form $x^{-s}$ plus a smooth function. At $x=b+$ only the second factor is singular. Compact subsets of $\Re s<1$, $\Re t<1$ therefore admit integrable dominating functions, including arbitrary spectral derivatives, because powers of $|\log x|$ remain integrable against $x^{-\sigma}$ for $\sigma<1$. This proves convergence and holomorphy. There is no additional condition $\Re(s+t)<1$ for a nonzero shift.

For $\Re s<0$, the Hurwitz Fourier expansion \cite[Eq.~25.11.9]{intdist:DLMFHurwitz} is absolutely convergent and its nonzero coefficients are
\begin{equation}\label{shj:fouriercoef}
\widehat\zeta_s(k)=\Gamma(1-s)(2\pi|k|)^{s-1}
\exp\!\left(-\frac{\pi\iu}{2}(1-s)\sgn k\right),\qquad k\ne0.
\end{equation}
The zero coefficient vanishes. For $\Re s,\Re t<0$, multiply the two absolutely convergent series and integrate. Orthogonality gives
\[
C(s,t;a)=\sum_{k\ne0}\widehat\zeta_s(k)\widehat\zeta_t(-k)e^{-2\pi\iu ka}.
\]
The positive and negative frequencies give exactly \eqref{shj:Cpoly}. Both sides are holomorphic on the larger connected domain stated in the proposition; the identity theorem extends the equality there.

To obtain \eqref{shj:Cbeta}, apply the same Fourier formula to $\zeta(1-w,a)$ and $\zeta(1-w,b)$, initially away from gamma poles. Euler reflection gives
$\Gamma(w)\Gamma(1-w)=\pi/\sin\pi w$ and
$1/\Gamma(t)=\Gamma(1-t)\sin\pi t/\pi$.
The coefficient of $\Li_w(z)$ reduces using
\[
\frac{\sin(\pi t)e^{-\pi\iu w/2}+\sin(\pi s)e^{\pi\iu w/2}}
{\sin(\pi w)}=e^{-\pi\iu(s-t)/2},\qquad w=2-s-t.
\]
The other coefficient is obtained by reversing the phase. This proves equality on an open set and hence meromorphically everywhere.
\end{proof}

\begin{corollary}[All ordinary spectral-jet correlations]\label{shj:ordinaryjets}
For integers $p,q\ge0$ and $\Re s,\Re t<1$,
\[
\int_0^1\zeta^{[p]}(s,x)\zeta^{[q]}(t,\{x+a\})\dd x
=\partial_s^p\partial_t^q C(s,t;a).
\]
Thus \eqref{shj:Cpoly} generates ordinary identities at nonpositive spectral integers, while continuation of \eqref{shj:Cbeta} generates finite-part identities at positive integers.
\end{corollary}
\begin{proof}
Differentiate under the integral using the compact domination established above. The formulas on the right are finite applications of the product and chain rules.
\end{proof}
The distinction between nonzero and zero shifts is essential. At $a=0$, the singularities coincide and the original convergence region changes to include a constraint on $s+t$. One must not substitute $a=0$ directly into a separated-singularity finite-part formula. The continuous ordinary log-gamma correlation will have a legitimate limit at zero, but that is a different assertion.


\begin{proposition}[Coefficient extraction at the double pole]\label{shj:Jcoefficient}
For $a\in(0,1)$,
\begin{equation}\label{shj:JfromC}
J_{mn}(a)=(-1)^{m+n}m!n![u^m v^n]C(1+u,1+v;a).
\end{equation}
The brackets mean the regular Laurent coefficient at the specified bidegree. In particular, the double-pole coefficient is $-1/(uv)$, not $1/(uv)$.
\end{proposition}
\begin{proof}
Write $Z_s=Z(s)$ for the periodic family of
Lemma~\ref{stconv:lem:Laurentdist}. Its raising-parameter expansion is
\begin{equation}\label{shj:Zlaurent}
 Z_{1+u}=-\stconvId/u+\sum_{m\ge0}(-1)^mG_m u^m/m!.
\end{equation}
The singular supports of $Z_{1+u}(x)$ and $Z_{1+v}(x+a)$ are $0$ and $b$, which are distinct. Choose disjoint neighborhoods of these points and a smooth partition of unity. On each neighborhood only one factor is a distribution and the other is smooth, so their product and integral are defined without any product of distributions at a common singularity. This gives the meromorphic continuation of the original integral.

The product of the regular coefficients in \eqref{shj:Zlaurent}, integrated with this partition, is exactly \eqref{shj:Jcutoff}. The pole terms give
\begin{align}
C(1+u,1+v;a)={}&-\frac1{uv}-\frac{R(v,a)}u-\frac{R(u,b)}v\notag\\
&+\sum_{m,n\ge0}\frac{(-1)^{m+n}}{m!n!}J_{mn}(a)u^mv^n.
\label{shj:fullLaurent}
\end{align}
The integrated product of $1-\delta_0$ and $1-\delta_b$ is $1-1-1=-1$ away from $a=0$, and each $G_m$ has mean zero. This proves the statement.
\end{proof}

\begin{corollary}[Polylogarithmic generator]\label{shj:polygen}
Set $E(u)=\Gamma(1-u)(2\pi)^u$ and
\begin{align}
\mathcal H(u,v;a)=E(u)E(v)\big[&e^{-\pi\iu(u-v)/2}\Li_{-u-v}(z)\notag\\
&+e^{\pi\iu(u-v)/2}\Li_{-u-v}(z^{-1})\big].
\label{shj:H}
\end{align}
Then
\begin{equation}\label{shj:Jpoly}
J_{mn}(a)=(-1)^{m+n}m!n![u^{m+1}v^{n+1}]\mathcal H(u,v;a).
\end{equation}
\end{corollary}
\begin{proof}
In \eqref{shj:Cpoly}, set $s=1+u$, $t=1+v$ and use $\Gamma(-u)=-\Gamma(1-u)/u$. Thus $C=\mathcal H/(uv)$, and \eqref{shj:JfromC} applies.
\end{proof}
This representation explains how polylogarithm \emph{order derivatives}, rather than only argument derivatives, enter the exact Stieltjes product calculus.

\section{An all-order linear Stieltjes closure theorem}\label{shj:closure}
The beta representation is particularly effective near $(s,t)=(1,1)$. Define
\begin{equation}\label{shj:AandT}
A(u,v)=\frac{\Gamma(1-u)\Gamma(1+u+v)}{\Gamma(1+v)},
\qquad T(u,v)=\frac{A(u,v)-1}{u(u+v)}.
\end{equation}
The apparent singularities in $T$ are removable: $A(0,v)=1$ and $A(u,-u)=1$. Its Taylor coefficients belong to the polynomial algebra generated by positive integer zeta values. More explicitly,
\begin{equation}\label{shj:logA}
\log A(u,v)=\sum_{k\ge2}\frac{\zeta(k)}k
\left[u^k+(-1)^k\big((u+v)^k-v^k\big)\right].
\end{equation}
Euler's constant cancels identically. This series converges in a neighborhood of the origin; it can also be used as an exact formal power series.

\begin{theorem}[All-index finite-part closure]\label{shj:closurethm}
Let $m,n\ge0$ and $d=m+n$. Define the analytic germ
\begin{align}
\mathcal F(u,v;a)={}&T(u,v)+T(v,u)\notag\\
&+(u+v)\big[T(u,v)R(u+v,a)+T(v,u)R(u+v,b)\big].
\label{shj:Fclosure}
\end{align}
Then
\begin{equation}\label{shj:closureformula}
\boxed{\displaystyle
J_{mn}(a)=\frac{\gamma_{d+1}(a)}{m+1}
+\frac{\gamma_{d+1}(b)}{n+1}
-(-1)^d m!n![u^m v^n]\mathcal F(u,v;a).}
\end{equation}
Every coefficient is obtained by a finite calculation. Apart from the two leading terms displayed explicitly, the right side contains only $\gamma_j(a)$ and $\gamma_j(b)$ with $0\le j\le d-1$, and a constant term. Its coefficients lie in
$\Q[\zeta(2),\ldots,\zeta(d+2)]$.
There are no products of generalized Stieltjes functions and no term involving $\gamma_d$.
\end{theorem}
\begin{proof}
Since $B(-u,1+u+v)=-A(u,v)/u$, equation \eqref{shj:Cbeta} becomes
\begin{equation}\label{shj:Cnear}
C(1+u,1+v;a)=-\frac{A(u,v)}u\zeta(1+u+v,a)
-\frac{A(v,u)}v\zeta(1+u+v,b).
\end{equation}
Substitute $A(u,v)=1+u(u+v)T(u,v)$ and the Laurent expansion \eqref{shj:laurent}. The result is
\begin{equation}\label{shj:Csplit}
C(1+u,1+v;a)=-\frac1{uv}
-\frac{R(u+v,a)}u-\frac{R(u+v,b)}v-\mathcal F(u,v;a).
\end{equation}
Here the apparent diagonal pole at $u+v=0$ has canceled exactly.

The regular coefficient of $-R(u+v,a)/u$ at $u^mv^n$ is
\[
-\frac{(-1)^{d+1}\gamma_{d+1}(a)}{(d+1)!}
\binom{d+1}{m+1}
=\frac{(-1)^d\gamma_{d+1}(a)}{(m+1)!n!}.
\]
The other term similarly contributes
$(-1)^d\gamma_{d+1}(b)/(m!(n+1)!)$.
Now apply \eqref{shj:JfromC} to get \eqref{shj:closureformula}.

To obtain a coefficient of total degree $d$ in $T$, only terms of total degree at most $d+2$ in $A$ are required. Formula \eqref{shj:logA} proves the asserted zeta coefficient algebra. In the product $(u+v)TR$, a Stieltjes index $j$ occupies degree $j+1$ before multiplication by $T$, so $j\le d-1$. All appearances of $R$ are linear. These observations prove the remaining claims.
\end{proof}

\begin{corollary}[Formal weight conservation]\label{shj:weight}
Assign weight $j+1$ to $\gamma_j(a)$ and $\gamma_j(b)$ and weight $k$ to $\zeta(k)$. Every monomial in \eqref{shj:closureformula} has weight $m+n+2$.
\end{corollary}
\begin{proof}
The degree-$k$ component of $\log A$ has coefficient weight $k$. Thus a coefficient of total degree $r$ in $T$ has weight $r+2$. A contribution from $(u+v)TR$ to total degree $d$ has $r+j+1=d$ and weight $(r+2)+(j+1)=d+2$. The leading Stieltjes terms have the same weight.
\end{proof}
This is a grading of the displayed symbolic algebra, not a theorem that the evaluated constants have no additional relations.

\subsection{The first four nontrivial identities}
The beginning of the universal kernel is
\begin{equation}\label{shj:Tstart}
T(u,v)=\zeta(2)-\zeta(3)v
+\frac{\zeta(4)+\zeta(2)^2}{2}(u^2+uv)
+\zeta(4)v^2+O((|u|+|v|)^3).
\end{equation}
Write $S_j(a)=\gamma_j(a)+\gamma_j(b)$. Coefficient extraction gives
\begin{align}
J_{00}(a)&=S_1(a)-2\zeta(2),\label{shj:J00}\\
J_{10}(a)&=\tfrac12\gamma_2(a)+\gamma_2(b)
+\zeta(2)S_0(a)-\zeta(3),\label{shj:J10}\\
J_{11}(a)&=\tfrac12 S_3(a)+2\zeta(2)S_1(a)+\zeta(3)S_0(a)
-\zeta(4)-\zeta(2)^2,\label{shj:J11}\\
J_{20}(a)&=\tfrac13\gamma_3(a)+\gamma_3(b)+2\zeta(2)S_1(a)
+2\zeta(3)\gamma_0(b)\notag\\
&\hspace{2cm}-3\zeta(4)-\zeta(2)^2.\label{shj:J20}
\end{align}
These are exact analytic identities under Definition~\ref{shj:defJ}, not conjectured PSLQ relations. The asymmetric coefficients in $J_{10}$ and $J_{20}$ are necessary; exchanging $a$ and $b$ exchanges the two Laurent indices.

Because $\gamma_0=-\psi$, the first identity takes the particularly simple form
\begin{equation}\label{shj:psipsi}
\FP\int_0^1\psi(x)\psi(\{x+a\})\dd x
=\gamma_1(a)+\gamma_1(b)-\frac{\pi^2}{3}.
\end{equation}
An equivalent identity containing only ordinary integrals is
\begin{align}
&\int_0^b\left[\psi(x)\psi(x+a)+\frac{\psi(a)}x\right]\dd x
+\int_0^a\left[\psi(b+t)\psi(t)+\frac{\psi(b)}t\right]\dd t\notag\\
&\qquad=\gamma_1(a)+\gamma_1(b)-\frac{\pi^2}{3}
+\psi(a)\log b+\psi(b)\log a.
\label{shj:psipsireg}
\end{align}
The endpoint singularities cancel inside each bracket. This version is useful when a conventional improper integral is preferable to distribution notation.

\subsection{A finite exact coefficient algorithm}
For target total index $d$, truncate \eqref{shj:logA} at degree $d+2$. Exponentiate by
\[
A=\sum_{j=0}^{\lfloor(d+2)/2\rfloor}\frac{(\log A)^j}{j!}
\pmod{\text{total degree}>d+2}.
\]
Exact polynomial division of $A-1$ by $u(u+v)$ has zero remainder. Let
$T(u,v)=\sum t_{ij}u^iv^j$. The correction coefficient needed in \eqref{shj:closureformula} is
\begin{align}
[u^mv^n]\mathcal F={}&t_{mn}+t_{nm}\notag\\
&+\sum_{\substack{i\le m,\ j\le n\\k=d-i-j-1\ge0}}
 t_{ij}\frac{(-1)^k}{k!}\binom{k+1}{m-i}\gamma_k(a)\notag\\
&+\sum_{\substack{j\le m,\ i\le n\\k=d-i-j-1\ge0}}
 t_{ij}\frac{(-1)^k}{k!}\binom{k+1}{m-j}\gamma_k(b).
\label{shj:coefficientalgorithm}
\end{align}
Only polynomial arithmetic over rational numbers and formal zeta symbols is used. The implementation does not perform floating-point relation recognition. Exact divisibility, reflection symmetry, linearity, and weight conservation provide independently checkable invariants of the generated formulas.


\providecommand{\scollCoeff}[1]{[#1]}
For the collision comparison, write the already proved linear closure as
\[
 J_{mn}(a)=e_{mn}+\sum_j(c_{mn,j}\gamma_j(a)+d_{mn,j}\gamma_j(b)).
\]
Use the local germs $\mathcal A(u,v)=A(v,u)$ and
$\mathcal B(u,v)=A(u,v)$ from \eqref{shj:AandT}, and put $w=u+v$.
The beta formula then reads
\begin{equation}\label{scoll:eq:kernel}
 K(u,v;a):=uv C(1+u,1+v;a)
 =-u\mathcal A(u,v)\zeta(1+w,b)-v\mathcal B(u,v)\zeta(1+w,a).
\end{equation}
Put $Q(u,v)=(u\mathcal A+v\mathcal B)/w$. Its numerator is divisible by $w$:
both germs equal one at $v=-u$. This removes the apparent diagonal
pole and supplies the finite coefficient construction used below.
\section{Coincident endpoints and canonical collision subtraction}
\label{scoll:sec:coincident}

A translated correlation has two separated logarithmic endpoint
divergences. At zero translation they merge into an additional
$x^{-2}$ singularity. We now show that the same coefficient table still
applies, but only after the merged singularity is subtracted explicitly.

For an integer $r\ge0$, put
\begin{equation}
 E_r(L)=\sum_{k=0}^r\frac{r!}{k!}L^k.
 \label{scoll:eq:E-polynomial}
\end{equation}
The identity $E_r-E_r'=L^r$ implies
\begin{equation}
 \int_\epsilon^1\frac{\log^r x}{x^2}\dd x
 =\frac{E_r(\log\epsilon)}{\epsilon}-r!.
 \label{scoll:eq:double-endpoint}
\end{equation}

\begin{definition}[Coincident-endpoint finite part]
\label{scoll:def:J}
Write $\gamma_j=\gamma_j(1)$ and define
\begin{align}
 J_{mn}^{\mathrm{coll}}=\lim_{\epsilon\downarrow0}\biggl\{
 &\int_\epsilon^1\gamma_m(x)\gamma_n(x)\dd x
       -\frac{E_{m+n}(\log\epsilon)}\epsilon\notag\\
 &+\frac{\gamma_n}{m+1}\log^{m+1}\epsilon
       +\frac{\gamma_m}{n+1}\log^{n+1}\epsilon\biggr\}.
 \label{scoll:eq:J-cutoff}
\end{align}
\end{definition}

\begin{theorem}[Coincident closure]
\label{scoll:thm:coincident}
The finite part \eqref{scoll:eq:J-cutoff} exists. With exactly the coefficients
of Theorem~\ref{shj:closurethm},
\begin{equation}
 J_{mn}^{\mathrm{coll}}=e_{mn}+\sum_{j=0}^{m+n+1}(c_{mn,j}+d_{mn,j})\gamma_j.
 \label{scoll:eq:J-closure}
\end{equation}
Equivalently it is the same mixed coefficient of
\begin{equation}
 K_0(u,v)=-(u\mathcal A(u,v)+v\mathcal B(u,v))\zeta(1+u+v).
 \label{scoll:eq:K0}
\end{equation}
In particular,
\begin{equation}
 J_{00}^{\mathrm{coll}}=2\gamma_1-\frac{\pi^2}{3}.
 \label{scoll:eq:J00}
\end{equation}
\end{theorem}
\begin{proof}
By \eqref{shj:localgamma}, the cutoff in Definition~\ref{scoll:def:J}
converges to the absolutely convergent expression
\begin{align}
 J_{mn}^{\mathrm{coll}}={}&-(m+n)!\notag\\
 &+\int_0^1\biggl\{
 \frac{\log^m x}{x}[\gamma_n(1+x)-\gamma_n]
 +\frac{\log^n x}{x}[\gamma_m(1+x)-\gamma_m]\notag\\
 &\hspace{35mm}+\gamma_m(1+x)\gamma_n(1+x)\biggr\}\dd x.
 \label{scoll:eq:J-convergent}
\end{align}
The two bracketed differences vanish linearly at zero, and
\eqref{scoll:eq:double-endpoint} supplies the constant $-(m+n)!$.

To identify the coefficients, start where the ordinary coincident
Hurwitz integral converges, for example $\Rea s,\Rea t<0$.
Fourier orthogonality gives
\[
 C_0(s,t)=2\Gamma(1-s)\Gamma(1-t)(2\pi)^{s+t-2}
       \cos\!\left(\frac{\pi(s-t)}2\right)\zeta(2-s-t).
\]
The Riemann functional equation, equivalently the specialization of the
Fourier calculation in Theorem~\ref{shj:masterthm}, gives
$uvC_0(1+u,1+v)=K_0(u,v)$. Meromorphic continuation is justified by
subtracting the local terms as above.

More explicitly, the holomorphic generating expression near zero is
\begin{align*}
 R_0(u,v)={}&-\frac1{1+u+v}\\
 &+\int_0^1\bigl\{
 x^{-1-u}[R(v,1+x)-R(v,1)]
 +x^{-1-v}[R(u,1+x)-R(u,1)]\\
 &\hspace{39mm}+R(u,1+x)R(v,1+x)\bigr\}\dd x.
\end{align*}
The first term is the analytic continuation of
$\int_0^1x^{-2-u-v}\dd x$, not a discarded divergence. The remaining
integrals converge locally uniformly near zero. On a starting convergence
region, and hence by continuation,
\[
 R_0=\frac{K_0+1+uR(u,1)+vR(v,1)}{uv}.
\]
Its Taylor coefficients are \eqref{scoll:eq:J-convergent}. Decomposing
\eqref{scoll:eq:K0} as $-Q-(u\mathcal A+v\mathcal B)R(u+v,1)$ now gives
\eqref{scoll:eq:J-closure} by the identical finite algebra used earlier.
\end{proof}

For the first case, an ordinary integral version is worth recording:
\begin{equation}
 \int_0^1\left\{\psi(1+x)^2
       -\frac{2[\psi(1+x)+\gamma]}x\right\}\dd x
 =1+2\gamma_1-\frac{\pi^2}{3}.
 \label{scoll:eq:digamma-square-convergent}
\end{equation}
The apparently singular numerator in the second term vanishes at zero.
There is no positivity contradiction in the negative value of $J_{00}^{\mathrm{coll}}$:
it is a finite part of a positive divergent square, not its ordinary
nonnegative integral.

\subsection{An exact convergent expansion at collision}
The relationship between the separated and coincident finite parts can
be expressed by an equality of convergent series, not merely an
asymptotic prescription.

\begin{theorem}[Collision polynomial and convergent regular part]
\label{scoll:thm:collision}
For $m,n\ge0$ define a polynomial
\begin{equation}
 \mathcal P_{mn}(L)=(-1)^{m+n+1}m!n!
       \scollCoeff{u^{m+1}v^n}\mathcal B(u,v)e^{-(u+v)L}.
 \label{scoll:eq:collision-polynomial}
\end{equation}
It has degree $m+n+1$, leading coefficient $1/(m+1)$, and coefficients
in positive integer zeta values. There is an exact expansion
\begin{equation}
 J_{mn}(a)=\frac{\mathcal P_{mn}(\log a)}a
                +J_{mn}^{\mathrm{coll}}+\sum_{k\ge1}h_{mn,k}a^k,
 \qquad 0<a<1,
 \label{scoll:eq:collision-series}
\end{equation}
whose regular power series converges locally uniformly for $|a|<1$.
Its coefficients are
\begin{align}
 h_{mn,k}={}&\frac{(-1)^{m+n+1}m!n!}{k!}
 \scollCoeff{u^{m+1}v^{n+1}}\notag\\
 &\quad\bigl(u\mathcal A(u,v)+(-1)^kv\mathcal B(u,v)\bigr)
               (1+u+v)_k\zeta(1+u+v+k).
 \label{scoll:eq:collision-coefficients}
\end{align}
For $m=n$, all odd-indexed $h_{mm,k}$ vanish. In particular,
\begin{equation}
 \lim_{a\downarrow0}\left\{J_{mn}(a)
             -\frac{\mathcal P_{mn}(\log a)}a\right\}=J_{mn}^{\mathrm{coll}}.
 \label{scoll:eq:collision-limit}
\end{equation}
\end{theorem}
\begin{proof}
The classical parameter Taylor series at one gives, locally uniformly
in the spectral variable,
\[
 \zeta(1+w,1-a)=\sum_{k\ge0}\frac{(1+w)_k}{k!}
                      \zeta(1+w+k)a^k,
\]
while the shift law gives
\[
 \zeta(1+w,a)=a^{-1-w}
        +\sum_{k\ge0}\frac{(1+w)_k}{k!}\zeta(1+w+k)(-a)^k.
\]
These are meromorphic in $w$ at zero, with the $k=0$ pole explicitly
visible. Substitution in \eqref{scoll:eq:kernel} yields
\begin{align}
 K(u,v;a)={}&-v\mathcal B(u,v)a^{-1-w}\notag\\
 &-\sum_{k\ge0}\frac{u\mathcal A(u,v)+(-1)^kv\mathcal B(u,v)}{k!}
                      (1+w)_k\zeta(1+w+k)a^k.
 \label{scoll:eq:K-collision}
\end{align}
The $k=0$ summand is exactly $K_0$; its apparent pole at $w=0$ is
removable. The remaining series is holomorphic for small $u,v$ and
$|a|<1$, uniformly on compact subsets. For example, on a fixed small
spectral polydisc its terms and spectral derivatives grow at most
polynomially in $k$ times powers of $\log k$, whereas $|a|^k$ decays
geometrically on compact subdiscs. Coefficient extraction is therefore
legitimate. The singular first term gives
\eqref{scoll:eq:collision-polynomial}, the constant gives
Theorem~\ref{scoll:thm:coincident}, and the other terms give
\eqref{scoll:eq:collision-coefficients}.

The highest power of $L$ comes from $B_0=1$ and has degree $m+n+1$;
a binomial coefficient gives its stated leading coefficient.
For $m=n$ and odd $k$, the factor
$u\mathcal A-v\mathcal B$ is antisymmetric under $u\leftrightarrow v$, while every other
factor in \eqref{scoll:eq:collision-coefficients} is symmetric. Its diagonal
mixed coefficient is zero. Finally the regular series has zero limit
at $a=0$, proving \eqref{scoll:eq:collision-limit}.
\end{proof}

The first collision polynomials are
\begin{equation}
 \mathcal P_{00}(L)=L,\qquad
 \mathcal P_{10}(L)=\frac{L^2}{2}+\zeta(2),\qquad
 \mathcal P_{01}(L)=L^2+\zeta(2).
 \label{scoll:eq:collision-examples}
\end{equation}
The asymmetry of the last two is required: the factor that becomes
singular at $a=0$ has changed its index. In particular,
\[
 J_{00}(a)-\frac{\log a}{a}\longrightarrow
                   2\gamma_1-\frac{\pi^2}{3}.
\]
There is no arbitrary additional finite constant in this limit. It is
fixed by the initial endpoint subtraction and agrees with the coincident
prescription.


\section{Ordinary integrals of all antiderivative orders}\label{shj:primitives}
The finite-part identities have a convergent counterpart. We first choose mean-zero periodic primitives, rather than implicitly choosing constants of integration at a singular endpoint.

\begin{definition}[Normalized primitive tower]\label{shj:Udef}
For integers $m\ge0$ and $r\ge1$, put
\begin{align}
\mathfrak h_j^{(r-1)}&=[u^j]\prod_{\nu=1}^{r-1}(1-u/\nu)^{-1},\notag\\
U_{m,r}(x)&=\frac{(-1)^{m+1}m!}{(r-1)!}
\sum_{j=0}^{m+1}\mathfrak h_{m+1-j}^{(r-1)}
\frac{\zeta^{[j]}(1-r,x)}{j!}.
\label{shj:Uformula}
\end{align}
The empty product convention gives $\mathfrak h_0^{(0)}=1$ and $\mathfrak h_j^{(0)}=0$ for $j>0$. Write $P_m=U_{m,1}$ and $Q_m=U_{m,2}$. Explicitly,
\begin{align}
P_m(x)&=\frac{(-1)^{m+1}}{m+1}\zeta^{[m+1]}(0,x),\label{shj:P}\\
Q_m(x)&=(-1)^{m+1}m!\sum_{j=0}^{m+1}\frac{\zeta^{[j]}(-1,x)}{j!}.
\label{shj:Q}
\end{align}
\end{definition}
These are normalized Hurwitz-jet primitives. The repository's primitive based at $1$ is $P_m(x)-P_m(1)$; it is not $P_m(x)$ itself. Likewise, zero-based negative polygamma conventions must be converted explicitly before using the following correlations.

\begin{lemma}[Primitive identities and regularity]\label{shj:Ulemma}
Each $U_{m,r}$ has mean zero on $(0,1)$. One has
\[
DP_m=G_m,\qquad DU_{m,r}=U_{m,r-1}\quad(r\ge2),\qquad D^rU_{m,r}=G_m
\]
as periodic distributions. The functions $P_m$ lie in $L^2(\stconvT)$, and for $r\ge2$, $U_{m,r}$ is a periodic $C^{r-2}$ function.
\end{lemma}
\begin{proof}
Consider the spectral primitive
\begin{equation}\label{shj:primitivegenerator}
V_r(u,x)=\frac{(-1)^r\zeta(1+u-r,x)}{(1+u-r)_r}
=-\frac{\zeta(1+u-r,x)}{(r-1)!u}
\prod_{\nu=1}^{r-1}(1-u/\nu)^{-1}.
\end{equation}
Its regular coefficient $(-1)^m m![u^m]V_r$ is \eqref{shj:Uformula}. Pointwise differentiation gives $\partial_xV_r=V_{r-1}$ for $r\ge2$, and $\partial_xV_1=\zeta(1+u,x)$. Comparing regular coefficients proves the pointwise primitive relations.

For $r=1$, the shift identity gives
\[
P_m(x)=\frac{\log^{m+1}x}{m+1}+P_m(1+x).
\]
This proves square integrability. Integration by parts against a periodic test function, with a lower cutoff, yields the boundary term $\phi(0)\log^{m+1}\epsilon/(m+1)$ and hence $DP_m=G_m$ by \eqref{stconv:def:G}.
For $r\ge2$, the singular term in $\zeta^{[j]}(1-r,x)$ is $x^{r-1}(-\log x)^j$, which tends to zero at $x=0$. The endpoint values at $0+$ and $1$ consequently match. Iterating the pointwise derivative relations shows periodic $C^{r-2}$ regularity and no extra jump terms. This proves the distributional relations. Finally, $\int_0^1\zeta(s,x)\dd x=0$ for $\Re s<1$, and its spectral derivatives also integrate to zero. Apply this to each term in \eqref{shj:Uformula}.
\end{proof}

\begin{theorem}[All-order convergent lifting]\label{shj:alllift}
Suppose the closure formula for fixed $m,n$ is written as
\begin{equation}\label{shj:linearform}
J_{mn}(a)=\sum_j c_j\gamma_j(a)+\sum_j d_j\gamma_j(b)+c,
\end{equation}
where all sums are finite and the coefficients are independent of $a$.
For $r,k\ge1$, $N=r+k$, define the ordinary integral
\[
W_{m,n}^{r,k}(a)=\int_0^1U_{m,r}(x)U_{n,k}(\{x+a\})\dd x.
\]
Then
\begin{align}
W_{m,n}^{r,k}(a)={}&(-1)^r\sum_j c_jU_{j,N}(a)
+(-1)^k\sum_j d_jU_{j,N}(b)\notag\\
&+\frac{(-1)^r c}{N!}B_N(a),
\label{shj:allliftformula}
\end{align}
where $B_N$ is the ordinary Bernoulli polynomial. Thus every such integral reduces to finitely many Hurwitz-zeta jets at the single spectral integer $1-r-k$ and zeta-valued coefficients. The highest required spectral derivative is $m+n+2$.
\end{theorem}
\begin{proof}
The integral exists by the $L^2$ property and Cauchy--Schwarz. Its mean in $a$ is zero by Fubini and the mean-zero normalization of both factors. In the distributional sense, periodic integration by parts gives
\begin{equation}\label{shj:Wderivative}
\frac{d^N}{da^N}W_{m,n}^{r,k}(a)=(-1)^rJ_{mn}(a),\qquad 0<a<1.
\end{equation}
The right side is smooth there.

We need enough endpoint regularity to fix all integration constants. From \eqref{shj:fouriercoef} and \eqref{shj:Zlaurent}, for $q\ne0$,
\begin{align}
\widehat G_m(q)=(-1)^{m+1}m![u^{m+1}]
\exp\bigg[&\left(\kappa+\log|q|+\frac{\pi\iu}{2}\sgn q\right)u\notag\\
&+\sum_{j\ge2}\frac{\zeta(j)}j u^j\bigg].
\label{shj:Gfourier}
\end{align}
In particular $\widehat G_m(q)=O_m((1+\log|q|)^{m+1})$.
The primitive relations imply
$\widehat U_{m,r}(q)=\widehat G_m(q)/(2\pi\iu q)^r$.
The Fourier coefficients of $W$ are therefore
$O((1+\log|q|)^{m+n+2}/|q|^N)$.
After $N-2$ derivatives the series is still absolutely convergent. Hence $W$ is periodic $C^{N-2}$.

Let $F(a)$ be the proposed right side of \eqref{shj:allliftformula}. It also has mean zero and is periodic $C^{N-2}$, using Lemma~\ref{shj:Ulemma} and the endpoint identities for Bernoulli polynomials. Its $N$th derivative on $(0,1)$ is $(-1)^rJ_{mn}$: the reflected primitive contributes $(-1)^N$, and $B_N^{(N)}=N!$. Thus $W-F$ is a polynomial of degree at most $N-1$ on $(0,1)$. Matching its endpoint derivatives of orders $0,\ldots,N-2$ forces every nonconstant coefficient to vanish, starting with the leading one. Its zero mean then forces the remaining constant to vanish.

The closure uses indices at most $m+n+1$, and $U_{j,N}$ uses spectral derivatives at most $j+1$. This gives the final derivative-order assertion.
\end{proof}

\subsection{The first-primitive correlation}
For $K_{mn}(a)=W_{m,n}^{1,1}(a)$, the theorem simplifies to
\begin{equation}\label{shj:Klift}
K_{mn}(a)=-\sum_j c_jQ_j(a)-\sum_j d_jQ_j(b)-\frac c2B_2(a),
\qquad K_{mn}''(a)=-J_{mn}(a).
\end{equation}
Equivalently, ordinary differentiation of the master integral at $(s,t)=(0,0)$ gives
\begin{equation}\label{shj:KfromC}
K_{mn}(a)=(-1)^{m+n}m!n![s^{m+1}t^{n+1}]C(s,t;a).
\end{equation}
The signs and factorials follow from \eqref{shj:P}; no finite-part operation occurs in this integral.

For example, \eqref{shj:J11} gives the fully explicit convergent identity
\begin{align}
K_{11}(a)={}&-\tfrac12\big(Q_3(a)+Q_3(b)\big)
-2\zeta(2)\big(Q_1(a)+Q_1(b)\big)\notag\\
&-\zeta(3)\big(Q_0(a)+Q_0(b)\big)
+\tfrac12\big(\zeta(4)+\zeta(2)^2\big)B_2(a).
\label{shj:K11}
\end{align}
This evaluates a product of two second spectral derivatives at zero using derivatives at $-1$ of order at most four.

An example with unequal integration orders is
\begin{equation}\label{shj:W21}
\int_0^1Q_0(x)P_0(\{x+a\})\dd x
=U_{1,3}(a)-U_{1,3}(b)-\frac{\zeta(2)}3B_3(a),
\end{equation}
where
\[
U_{1,3}(x)=\frac78\zeta(-2,x)+\frac34\zeta^{[1]}(-2,x)
+\frac14\zeta^{[2]}(-2,x).
\]
Thus the same finite closure gives exact identities involving successive antiderivatives rather than merely a formal recurrence for integration constants.

\section{Explicit log-gamma and mixed identities}\label{shj:loggamma}
Let
\[
f(x)=\log\Gamma(x)-\frac\ell2,\qquad
\mathcal K(a)=\int_0^1 f(x)f(\{x+a\})\dd x.
\]
Lerch's identity gives $P_0=-f$, so $\mathcal K=K_{00}$. No regularization is used in this definition: logarithmic singularities are square integrable.

\begin{theorem}[Translated log-gamma correlation]\label{shj:gammathm}
For $0<a<1$,
\begin{align}
\mathcal K(a)={}&(1+\zeta(2))B_2(a)
-\zeta^{[1]}(-1,a)-\zeta^{[1]}(-1,b)\notag\\
&-\frac12\big[\zeta^{[2]}(-1,a)+\zeta^{[2]}(-1,b)\big].
\label{shj:gammahurwitz}
\end{align}
Equivalently, for $z=e^{2\pi\iu a}$,
\begin{equation}\label{shj:gammapoly}
\mathcal K(a)=\frac1{2\pi^2}\Re\left[
\Li_2^{[2]}(z)-2\kappa\Li_2^{[1]}(z)
+\left(\kappa^2+\frac{\pi^2}{4}\right)\Li_2(z)\right].
\end{equation}
For the uncentered product,
\begin{equation}\label{shj:gammauncentered}
\int_0^1\log\Gamma(x)\log\Gamma(\{x+a\})\dd x
=\frac{\ell^2}{4}+\mathcal K(a).
\end{equation}
Finally,
\begin{equation}\label{shj:gammasecond}
\mathcal K''(a)=\frac{\pi^2}{3}-\gamma_1(a)-\gamma_1(b).
\end{equation}
\end{theorem}
\begin{proof}
Apply \eqref{shj:Klift} to $J_{00}=\gamma_1(a)+\gamma_1(b)-2\zeta(2)$. Since
\[
Q_1(x)=\zeta(-1,x)+\zeta^{[1]}(-1,x)+\tfrac12\zeta^{[2]}(-1,x)
\]
and $\zeta(-1,x)=-B_2(x)/2$, the result is \eqref{shj:gammahurwitz}.

For an independent Fourier derivation of \eqref{shj:gammapoly}, differentiate \eqref{shj:fouriercoef} at $s=0$. The positive-frequency coefficient of $f$ is
\[
\widehat f(n)=\frac1{4n}-\frac{\iu(\kappa+\log n)}{2\pi n},\qquad n\ge1.
\]
The negative-frequency coefficient is its conjugate. These coefficients also give the classical Kummer Fourier expansion. The $L^2$ correlation therefore equals the absolutely convergent series
\begin{equation}\label{shj:gammafourier}
\mathcal K(a)=\frac1{2\pi^2}\sum_{n\ge1}
\frac{\cos(2\pi na)}{n^2}
\left[(\kappa+\log n)^2+\frac{\pi^2}{4}\right].
\end{equation}
Taking the first two spectral derivatives of the absolutely convergent polylogarithm series at order $2$ proves \eqref{shj:gammapoly}. The mean of $f$ is zero, so expanding the uncentered product proves \eqref{shj:gammauncentered}. Equation \eqref{shj:gammasecond} is the $m=n=0$ case of \eqref{shj:Klift}.
\end{proof}

\begin{corollary}[A mixed log-gamma--digamma identity]\label{shj:mixedgamma}
The same finite-part convention gives
\begin{align}
\FP\int_0^1\log\Gamma(x)\psi(\{x+a\})\dd x
={}&\frac12\big[\zeta^{[2]}(0,b)-\zeta^{[2]}(0,a)\big]\notag\\
&+\zeta(2)(2a-1).
\label{shj:mixedformula}
\end{align}
It is unchanged when $\log\Gamma(x)$ is replaced by $f(x)$.
\end{corollary}
\begin{proof}
Since $DP_0=G_0$ with no contact term at this first primitive stage,
$\mathcal K'(a)=\FP\int f(x)\psi(\{x+a\})\dd x$.
Differentiate \eqref{shj:gammahurwitz} using
\[
\partial_x\zeta^{[j]}(-1,x)
=\zeta^{[j]}(0,x)-j\zeta^{[j-1]}(0,x).
\]
The first spectral derivatives at zero cancel, and the terms from
$\zeta(0,x)=1/2-x$ simplify to $\zeta(2)(2a-1)$. This gives \eqref{shj:mixedformula}. The canonical periodic finite part of $\psi$ has mean zero, so subtracting $\ell/2$ from the log-gamma factor does not change the answer.
\end{proof}
For example, its left side with the centered function equals the entirely convergent expression
\begin{align*}
&\int_0^b f(x)\psi(x+a)\dd x\\
&\quad+\int_0^a\left[(f(b+t)-f(b))\psi(t)+f(b)\psi(1+t)\right]\dd t
-f(b)\log a.
\end{align*}
The recurrence $\psi(1+t)=\psi(t)+1/t$ removes the endpoint pole explicitly. This form also provides a direct numerical check independent of differentiating a finite-part integral.

\subsection{Half shift and the zero-shift limit}
At $a=1/2$, use $\Li_s(-1)=(2^{1-s}-1)\zeta(s)$. With $L_2=\log2$,
\begin{align}
\mathcal K(1/2)={}&-\frac{\zeta^{[2]}(2)}{4\pi^2}
+\frac{\kappa-L_2}{2\pi^2}\zeta^{[1]}(2)\notag\\
&+\frac{L_2^2+2\kappa L_2-\kappa^2-\pi^2/4}{24}.
\label{shj:gammahalf}
\end{align}
This formula concerns a \emph{translated} correlation, not the reflected product $f(x)f(1-x)$.

The absolutely convergent series \eqref{shj:gammafourier} is continuous at $a=0$. Its value there recovers the classical second log-gamma moment \cite{stieltjes:EspinosaMoll1}:
\begin{align}
\int_0^1\log^2\Gamma(x)\dd x={}&
\frac{\ell^2}{4}+\frac{\zeta^{[2]}(2)}{2\pi^2}
-\frac{\kappa\zeta^{[1]}(2)}{\pi^2}
+\frac{\kappa^2}{12}+\frac{\pi^2}{48}.
\label{shj:classicalmoment}
\end{align}
Equivalently, taking the continuous limit in \eqref{shj:gammahurwitz} gives
\[
\mathcal K(0)=\frac{1+\zeta(2)}6-2\zeta^{[1]}(-1)-\zeta^{[2]}(-1).
\]
These are useful consistency checks, not priority claims. In contrast,
$J_{00}(a)=\log a/a+O(1)$ as $a\downarrow0$, by \eqref{shj:localgamma}. The finite-part product has no continuous zero-shift specialization. The ordinary correlation remains continuous but its second derivative develops the expected singularity.

\section{Rational shifts and exact arithmetic coordinates}\label{shj:rationals}
For integers $q\ge2$ and $1\le p<q$, residue-class decomposition gives the classical finite Fourier identity
\begin{equation}\label{shj:DFT}
\Li_s(e^{2\pi\iu p/q})
=q^{-s}\sum_{r=1}^q e^{2\pi\iu pr/q}\zeta(s,r/q).
\end{equation}
It is proved by absolutely convergent series for $\Re s>1$ and then by analytic continuation. Spectral differentiation yields
\begin{align}
\Li_s^{[j]}(e^{2\pi\iu p/q})
=q^{-s}\sum_{r=1}^q e^{2\pi\iu pr/q}
\sum_{h=0}^j\binom jh(-\log q)^{j-h}\zeta^{[h]}(s,r/q).
\label{shj:DFTjets}
\end{align}
At $s=1$ the poles of the finite sum cancel because the roots sum to zero. One must combine the pole terms before evaluating; separately substituting $\zeta(1,r/q)$ is invalid.

At a rational shift, Theorem~\ref{shj:closurethm} reduces all Stieltjes correlations to rational-argument Stieltjes constants and ordinary zeta values. Theorem~\ref{shj:alllift} places every primitive correlation at a negative integer spectral argument. Equation \eqref{shj:DFTjets} instead evaluates the polylogarithmic side using positive-order Hurwitz jets. The two routes give independent choices of arithmetic coordinates for numerical replay and exact symbolic comparison.

For a Dirichlet character modulo $q$, including an imprimitive one with its actual modulus, the identity
\[
\sum_{r=1}^q\chi(r)\zeta(s,r/q)=q^sL(s,\chi)
\]
and its spectral derivatives transform the character components into Dirichlet-$L$ jets. This gives a reduction framework, not a claim of numerical independence or minimality of the resulting coordinates.

\begin{corollary}[Small-denominator digamma correlations]\label{shj:rationalpsi}
Write $\gamma_1=\gamma_1(1)$. Then
\begin{align}
\Pi_{0,0}(1/2)&=2\gamma_1-4\gamma\log2-2\log^22-\frac{\pi^2}{3},\label{shj:rathalf}\\
\Pi_{0,0}(1/3)&=2\gamma_1-3\gamma\log3-\frac32\log^23-\frac{\pi^2}{3},\label{shj:ratthird}\\
\Pi_{0,0}(1/4)&=2\gamma_1-6\gamma\log2-7\log^22-\frac{\pi^2}{3}.
\label{shj:ratquarter}
\end{align}
\end{corollary}
\begin{proof}
Expand at $s=1$ the exact identities
\begin{align*}
\zeta(s,1/2)&=(2^s-1)\zeta(s),\\
\zeta(s,1/3)+\zeta(s,2/3)&=(3^s-1)\zeta(s),\\
\zeta(s,1/4)+\zeta(s,3/4)&=(4^s-2^s)\zeta(s).
\end{align*}
Compare the coefficient of $s-1$ with \eqref{shj:laurent} and insert the resulting sums into \eqref{shj:psipsi}. In the first line the correlation contains two copies of $\gamma_1(1/2)$, which accounts for its different coefficients.
\end{proof}
Higher indices require no new search: expand these finite exponential factors to the needed order and use \eqref{shj:closureformula}. At larger denominators, retaining a nontrivial character block is legitimate; failing to reduce that block does not establish transcendence or independence.


The canonical finite-part and primitive identities above are audited
against the independent coefficient engines and their separately labeled
numerical replays under \src{verification/eighth-replay/}. The full contact
law is proved in the next chapter before derivatives of these finite-part
identities are used. The singular same-point finite part and the continuous
ordinary log-Gamma autocorrelation have different zero-shift behavior.
